\chapter{总结}

本文围绕复杂环境下GNSS/INS松耦合组合导航中观测噪声的非高斯性与重尾特性，提出并实现了一种基于自适应最大相关熵卡尔曼滤波（Adaptive MCKF）的鲁棒松耦合方案，并完成了相应的算法工程化实现。本文的主要贡献与结论如下：

\begin{enumerate}
\item \textbf{基于误差状态的松耦合滤波框架}：构建了包含位置、速度、姿态误差及传感器偏置的15维误差状态动力学模型，采用误差状态卡尔曼滤波（ESKF）架构将名义导航态与误差态解耦，为在工程中插入鲁棒滤波器提供了清晰可扩展的结构基础\cite{yan2019strapdown}。

\item \textbf{自适应最大相关熵滤波器的融入}：将最大相关熵准则（MCC）引入到松耦合滤波环节，用以替代传统的二阶矩准则，利用相关熵在核空间对高阶统计信息的敏感性提升了滤波器对突发观测异常和重尾噪声的鲁棒性。同时提出基于新息协方差匹配的自适应核带宽调整策略，使滤波过程能够在时变噪声环境中动态平衡预测与观测的影响，降低因异常观测导致的估计偏差\cite{izanloo2016kalman,xie2025adaptive}。

% \item \textbf{工程化与可复现的算法库}：实现了模块化、配置化的组合导航算法库，支持IMU高频预测与GNSS低频校正的双速率调度，误差传播采用矩阵指数离散化，传感器偏置以一阶马尔可夫过程建模。代码设计便于替换滤波器核心、调整核带宽策略并进行横向对比，便于在不同数据集与平台上复现与迁移。

\item \textbf{工程化处理策略}：在面向工程实现方面，提出并实现了若干实用处理措施，包括对新息协方差矩阵的正定化方案、双速率传感器的匹配与利用策略，以及对GNSS/INS误差模型采用的不完全闭环（部分闭环反馈）机制，以提高算法在含多径、信号衰减等复杂环境中的鲁棒性与稳定性。

\item \textbf{初步实验评估}：在包含多径效应与信号衰减的模拟城市峡谷场景中，三次重复实验表明，自适应MCKF能够稳定抑制误差发散，位置、速度和姿态RMSE分别达到 $4.24\pm0.25~\mathrm{m}$、$0.057\pm0.002~\mathrm{m/s}$ 和 $4.87\pm1.46^{\circ}$；相比之下，EKF 的位置与速度误差显著更大，且存在明显的长时漂移与局部发散现象。位置 RMSE 的统计检验结果为 $p=0.422$，说明在当前样本规模下差异尚未达到严格显著性水平，但从误差曲线、峰值误差与分位数统计来看，本文方法在鲁棒性与工程可用性方面更具优势。
\end{enumerate}

总体而言，本文提出的方法在方法论与工程实现上均具有一定的实用性：方法层面通过引入相关熵与自适应核带宽增强了对非高斯噪声的鲁棒性；实验层面则表明，该方法在复杂城市峡谷场景下能够有效抑制松耦合导航中的长时漂移和异常放大问题。后续工作可继续在更多典型工况下扩展对比实验，并进一步研究紧耦合结构与在线参数自适应策略，以提升实际部署能力。